\documentclass[14pt, a4paper, oneside]{extarticle}
\usepackage[T2A]{fontenc}
\usepackage[utf8]{inputenc}
\usepackage[russian]{babel}
\usepackage{hyperref}
\usepackage{xcolor}
\usepackage{geometry}
\geometry{top=1cm, bottom=0.5cm, left=1.25cm, right=1.25cm}
\usepackage{graphicx}
\usepackage{float}
\usepackage{wrapfig}
\usepackage{fontawesome5}
\usepackage{ragged2e}
\usepackage{microtype}
\usepackage{amsmath}
\usepackage{amssymb}

\hypersetup{
    colorlinks=black,
    allcolors=blue,
    pdfborderstyle={/S/U/W 0},
    allbordercolors=blue
}

\begin{document}

Рассматривается система линейных алгебраических уравнений
\begin{equation}\label{eq:slae}
A x = b,
\end{equation}
где $A \in \mathbb{R}^{n \times n}$~--- квадратная матрица с ненулевыми
диагональными элементами $a_{ii} \neq 0$,
 $x = (x_1, \dots, x_n)^{T}$~--- вектор неизвестных,
 $b = (b_1, \dots, b_n)^{T}$~--- вектор правой части.

Представим матрицу $A$ в виде суммы диагональной и внедиагональной частей:
\begin{equation}\label{eq:split}
A = D + L + U, \qquad
D = \operatorname{diag}(a_{11}, \dots, a_{nn}),
\end{equation}
где $D$~--- диагональная матрица, $L$ и $U$~--- строго нижняя и строго
верхняя треугольные части матрицы $A$. Тогда \eqref{eq:slae} записывается
как $D x = b - (L + U) x$, и итерационный процесс принимает матричный вид
\begin{equation}\label{eq:jacobi_matrix}
x^{(k+1)} = D^{-1}\bigl(b - (L + U)\, x^{(k)}\bigr)
          = \bigl(I - D^{-1} A\bigr)\, x^{(k)} + D^{-1} b,
\qquad k = 0, 1, 2, \dots
\end{equation}
где $x^{(0)}$~--- начальное приближение. В компонентной записи итерации
метода Якоби имеют вид
\begin{equation}\label{eq:jacobi_comp}
x_i^{(k+1)} = \frac{1}{a_{ii}}
\Bigl( b_i - \sum_{j \ne i} a_{ij}\, x_j^{(k)} \Bigr),
\qquad i = 1, \dots, n.
\end{equation}

\textbf{Теорема (достаточное условие сходимости).} Если матрица $A$ обладает строгим диагональным преобладанием,
\begin{equation}\label{eq:diagdom}
|a_{ii}| > \sum_{j \ne i} |a_{ij}|, \qquad i = 1, \dots, n,
\end{equation}
то метод Якоби сходится для любого $x^{(0)}$. 

Итерации прекращаются при достижении заданной точности $\varepsilon > 0$ по разности
соседних приближений:
\begin{equation}\label{eq:stop}
\frac{\bigl\| x^{(k)} - x^{(k-1)} \bigr\|}{\bigl\| x^{(k)} \bigr\|} \le \varepsilon.
\end{equation}

\end{document}